\section{A certified fixed-field branch at the coefficient prime 269}
\label{sec:cubic-fixed-269-classes}

This section concerns the fixed field
\(E=\mathbb Q(\omega,\sqrt[3]{269})\), not the growing hosts in
AM-H06--AM-H08.  Put \(p=269\), and let \(S\) comprise the primes
above \(p\) and the infinite places.  For a finite abelian
\(p\)-group \(C\), write
\(d_p(C)=\dim_{\mathbb F_p}(C/pC)\).  The strict group
\(S_p(E,e)\) is exactly the group defined in AC-K10.
The ordinary class-number certificate in AC-K13 is the only
computer-assisted input used below; no new number-field computation
is asserted.

\begin{theorem}[AC-K31: the fixed 269 branch and its exact finite class-group criterion]
\label{AC-K31}
For every integer \(e\geq1\), there is an isomorphism of finite groups
\[
 S_p(E,e)^\vee\simeq
 H^2(G_{E,S},\mathbb Z_p(e+1))
 \simeq K_{2e}(\mathcal O_E)_{(p)}.                 \tag{K31.1}
\]
In particular,
\[
 e\equiv0\pmod{268}
 \quad\Longrightarrow\quad
 S_{269}(E,e)=K_{2e}(\mathcal O_E)_{(269)}=0.       \tag{K31.2}
\]

There is also an exact criterion for every other residue class.
Let
\[
 L=E(\mu_p),\qquad
 \Delta=\operatorname{Gal}(L/E),\qquad
 \bar\omega:\Delta\longrightarrow\mathbb F_p^\times
\]
be the mod-\(p\) cyclotomic character.  Then
\(\Delta\simeq\mathbb F_p^\times\), \([L:\mathbb Q]=1608\),
and
\[
\begin{split}
 d_p\bigl(S_p(E,e)\bigr)
 &=d_p\bigl(K_{2e}(\mathcal O_E)_{(p)}\bigr)\\
 &=\dim_{\mathbb F_p}
       \bigl(\operatorname{Cl}(L)/p\bigr)^{\bar\omega^{-e}}.
\end{split}                                                   \tag{K31.3}
\]
The superscript denotes the eigenspace for the natural action on
ideal classes, with precisely the displayed inverse character.
Consequently all the groups in (K31.1), for all \(e\geq1\),
vanish if and only if \(\operatorname{Cl}(L)_{(269)}=0\).
This last class group has not been computed here.  Formula
(K31.3) gives a minimal-generator and vanishing criterion; it
does not give the orders of the nonzero groups.

Finally, every finite layer \(E_n\) of the cyclotomic
\(\mathbb Z_{269}\)-extension of \(E\) satisfies
\[
                   \operatorname{Cl}(E_n)_{(269)}=0.           \tag{K31.4}
\]
Neither theorem A nor an Iwasawa main conjecture is used.
\end{theorem}
\begin{proof}
\proofstep{1. The unique prime and the two cyclotomic images.}
The prime \(p\equiv2\pmod3\) is inert in \(K=\mathbb Q(\omega)\).
At its unique prime the polynomial \(X^3-p\) is Eisenstein.
Thus \(E\) has a unique prime \(v\) above \(p\), with
ramification index three and residue degree two over \(\mathbb Q_p\).
The local extension \(E_v/\mathbb Q_p\) is the splitting field
of \(X^3-p\): its unramified quadratic subfield contains
\(\mu_3\).  Its Galois group is \(S_3\), because unramified
Frobenius acts on \(\mu_3\) by inversion.  The only nontrivial
abelian quotient of \(S_3\) has order two.  Accordingly the
only nontrivial abelian subfield of \(E_v\) is this unramified
quadratic field.  Every \(\mathbb Q_p(\mu_{p^n})/\mathbb Q_p\)
is totally ramified.  Their intersection with \(E_v\) is
therefore \(\mathbb Q_p\).

Globally \(E/\mathbb Q\) is likewise the \(S_3\)-splitting
field of \(X^3-p\).  Its only nontrivial abelian subfield is
\(K\).  The latter ramifies at three, whereas every
\(\mathbb Q(\mu_{p^n})\) is unramified away from \(p\).
Hence
\[
 E\cap\mathbb Q(\mu_{p^\infty})=\mathbb Q,
 \qquad
 E_v\cap\mathbb Q_p(\mu_{p^\infty})=\mathbb Q_p.     \tag{K31.5}
\]
Both the global and the local cyclotomic images are consequently
the whole \(\mathbb Z_p^\times\).  For
\(A=\mathbb Q_p/\mathbb Z_p(-e)\), restriction thus identifies
\(H^0(E,A)\) with \(H^0(E_v,A)\), as the same subgroup of
\(A\).  There is only this one finite local term.  The last
map in (K10.1) is therefore an isomorphism, proving (K31.1).
In particular no local correction has been silently omitted.

\proofstep{2. A one-prime finite-coefficient calculation.}
For any totally imaginary number field \(F\) with exactly one
prime \(w\) above an odd prime \(p\), let \(S_F\) consist
exactly of that prime and the infinite places.  We claim a natural
isomorphism
\[
 H^2(G_{F,S_F},\mu_p)
       \simeq \operatorname{Cl}(\mathcal O_F[1/p])/p.          \tag{K31.6}
\]
Here is a direct duality proof, including the possible local
term.  The final map in finite Poitou--Tate duality for
\(\mu_p\) is
\[
 H^2(F_w,\mu_p)\longrightarrow H^0(F,\mathbb F_p)^\vee.
\]
Under local Tate duality it is dual to restriction of constant
\(\mathbb F_p\)-valued functions.  It is an isomorphism, since
there is exactly one \(w\).  Thus the preceding global
\(H^2\) is its strict localization kernel.  Poitou--Tate
duality identifies that kernel with the dual of
\[
 \ker\bigl(H^1(G_{F,S_F},\mathbb F_p)
                 \longrightarrow H^1(F_w,\mathbb F_p)\bigr).
\]
These are precisely the exponent-\(p\) abelian characters
unramified everywhere and split at \(w\).  Global class field
theory identifies them with
\(\operatorname{Hom}(\operatorname{Cl}(\mathcal O_F[1/p]),
\mathbb F_p)\).  Taking the finite dual proves (K31.6).
All the maps are natural under automorphisms of \(F\); the
result is equivariant for such automorphisms, with their usual
action on ideal classes.  No Tate twist is inserted on the
class group.

The exact duality inputs are Sharifi, \emph{Iwasawa Theory},
Appendix A, Theorems A.0.3, A.0.6 and A.0.8; the last theorem's
printed map construction makes the preceding restriction map
explicit:
\url{https://www.math.ucla.edu/~sharifi/notes/iwasawa-app-a.html}.
The source gives a proof sketch of the nine-term sequence; this
argument uses the stated duality theorem and does not claim a
new proof of that theorem.

\proofstep{3. Reduction of compact degree-two cohomology.}
For every such field \(F\) and \(e\geq1\), one has
\[
 H^2(G_{F,S_F},\mathbb Z_p(e+1))/p
       \simeq H^2(G_{F,S_F},\mathbb F_p(e+1)).                 \tag{K31.7}
\]
To justify the surjectivity, rather than assume it, use that
\(G_{F,S_F}\) has \(p\)-cohomological dimension at most two
and finite cohomology with finite \(p\)-primary coefficients.
These precise background facts are recorded in Lim--Sharifi,
\emph{Nekov\'a\v r duality over \(p\)-adic Lie extensions of
global fields}, Documenta Math. \textbf{18} (2013), 621--678,
\S4.1, first paragraph, printed p.~657:
\url{https://ems.press/content/serial-article-files/26190}.
Let \(T=\mathbb Z_p(e+1)\).  Continuous cochains with values
in \(T\) are the inverse limit of the cochains with values in
\(T/p^n\); the transition maps of cochains are surjective,
by lifting each of the finitely many values of a locally
constant cochain.  The inverse-limit cohomology exact sequence
has \(\varprojlim H^3(G_{F,S_F},T/p^n)=0\), and its
\(\varprojlim^1 H^2\) term is zero: each of these finite
cohomology groups satisfies the Mittag--Leffler condition in
the inverse system.  Hence \(H^3(G_{F,S_F},T)=0\).
The long exact sequence of
\(0\to T\xrightarrow{p}T\to T/p\to0\) now gives (K31.7).
The compact degree-two group is finite by AC-K10.  It is zero
if and only if its reduction modulo \(p\) is zero.

For \(e\equiv0\pmod{268}\), the coefficient
\(\mathbb F_p(e+1)\) is \(\mu_p\) already over \(E\).
AC-K13 certified the ordinary class group
\(\operatorname{Cl}(E)\simeq(\mathbb Z/7)^2\), by the
unconditional command \texttt{bnfcertify(b)} returning
\texttt{1}.  Its \(S\)-class quotient has zero \(p\)-part.
Equations (K31.6)--(K31.7) and (K31.1) prove (K31.2).

\proofstep{4. All finite cyclotomic characters.}
Equation (K31.5) gives \([L:E]=268\), and the same local
disjointness shows that the unique place over \(p\) has local
degree \(6\cdot268\).  Thus there is only one such place
\(\mathfrak P\) in \(L\).  The two tame ramification indices
three and 268 are coprime, so
\[
 e(\mathfrak P/p)=804,\quad f(\mathfrak P/p)=2,
 \qquad p\mathcal O_L=\mathfrak P^{804}.                       \tag{K31.8}
\]
In particular the order of \([\mathfrak P]\) divides 804,
which is prime to 269.  Quotienting the class group by that
class has no effect modulo \(p\).

The extension \(L/E\) is unramified outside \(S\) and has
degree prime to \(p\).  Restriction and the
Hochschild--Serre spectral sequence, whose positive-degree
\(\Delta\)-cohomology vanishes by averaging, give
\[
\begin{split}
 H^2(G_{E,S},\mathbb F_p(e+1))
 &\simeq H^2(G_{L,S_L},\mathbb F_p(e+1))^\Delta\\
 &\simeq
 \left(H^2(G_{L,S_L},\mu_p)\otimes_{\mathbb F_p}
                 \mathbb F_p(e)\right)^\Delta\\
 &\simeq
 \bigl(\operatorname{Cl}(L)/p\bigr)^{\bar\omega^{-e}}.
\end{split}                                                    \tag{K31.9}
\]
For the middle equality, \(\mathbb F_p(e)\) is trivial on
\(G_L\); its remaining \(\Delta\)-action is
\(\bar\omega^e\).  Together with the natural equivariance in
(K31.6), this proves the inverse-character convention in the
last line.  Combine (K31.7), (K31.9) and (K31.1) to obtain
(K31.3); a finite abelian group and its Pontryagin dual have
the same elementary-divisor ranks.

All the \(268\) characters of \(\Delta\) are defined over
\(\mathbb F_p\), and \(\mathbb F_p[\Delta]\) is semisimple
by averaging.  Positive integers \(e\) cover all their
residue classes.  Hence the vanishing of all the eigenspaces
in (K31.9) is equivalent to \(\operatorname{Cl}(L)/p=0\),
and thus to \(\operatorname{Cl}(L)_{(p)}=0\).  No assertion
about that uncomputed group is needed for the equivalence.

\proofstep{5. The pro-269 tower does not supply the missing characters.}
Let \(E_n\) be its degree \(p^n\) layer.  The local
disjointness in (K31.5) implies that \(E_{n+1}/E_n\) has
exactly one ramified finite prime, totally ramified of degree
\(p\).  No other finite prime ramifies and there are no real
places.  The certified base case is
\(\operatorname{Cl}(E)_{(p)}=0\).  Inductively, Chevalley's
ambiguous class-number formula gives
\[
 \#\operatorname{Cl}(E_{n+1})^{\operatorname{Gal}(E_{n+1}/E_n)}
 =\frac{\#\operatorname{Cl}(E_n)}
 {[\mathcal O_{E_n}^{\times}:
   \mathcal O_{E_n}^{\times}\cap
     N_{E_{n+1}/E_n}(E_{n+1}^{\times})]} .                    \tag{K31.10}
\]
Its \(p\)-primary part is therefore trivial.  A nonzero finite
abelian \(p\)-group with a cyclic order-\(p\) action has a
nonzero fixed element: on its nonzero \(p\)-torsion vector
space the operator \(\sigma-1\) is nilpotent, since
\((\sigma-1)^p=0\).  Thus the whole \(p\)-class group is
zero, proving (K31.4).  The precise formula, already used in
AC-K16, is C.-F.~Yu, \emph{Ambiguous class number formulas},
Theorem~1.1, equation~(1.2), printed p.~1, with its ordinary
interpretation on p.~2 and complete proof in \S2, pp.~2--3:
\url{https://archive.mpim-bonn.mpg.de/id/eprint/3296/1/preprint_2014_68.pdf}.
This pro-\(p\) tower does not adjoin the nontrivial
prime-to-\(p\) cyclotomic characters in \(\Delta\), so it
does not turn (K31.2) into all-weight vanishing.
\end{proof}

\paragraph{Scope of the arithmetic consequence.}
For the fixed field \(E\), the entire coefficient-269 term is
zero on the branch \(e\equiv0\pmod{268}\).  The literal
\texttt{GD.gen} condition at 269 excludes this branch, since
\(\bar\kappa^{2e}\) is then trivial on inertia.  Thus this
result supplies neither the remaining actual Type-II mixed
tangent vanishing at 269 nor a bound for the growing hosts of
AM-H06--AM-H08.  Those hosts can have several primes above 269
and nonzero class groups.  Even a certificate
\(\operatorname{Cl}(E(\mu_{269}))_{(269)}=0\) would address
the fixed field only.  Regularity of 269 over \(\mathbb Q\)
does not give that certificate for this nonabelian compositum.

\paragraph{Finite \'etale and motivic dictionary.}
On the arithmetic curve \(\operatorname{Spec}\mathcal O_E[1/p]\),
adjoining \(\mu_p\) is the finite \(\Delta\)-cover
\(\operatorname{Spec}\mathcal O_L[1/p]\).  Equation (K31.9)
is descent of the finite Tate local system, followed by its
character decomposition; (K31.6) is arithmetic duality and
class field theory.  The even K-group is identified by the
integral comparison already cited and scoped in AC-K10.
The previously audited SGA~1, Expos\'e~V, \S\S4 and 7,
supplies the finite-cover/Galois-category dictionary only.
It does not compute an unverified class group or bound a
nonzero specialization.
